% answers_short.tex —— 2014 年·广东卷·理科 · 简短答案速查
% 由 tools/make_short_answers.py 自 answers.tex 的 \gksolhead/\gkanswer 提取，
% 并与 2014/2014_数学真题_广东理_答案要点.md 逐题比对后生成，勿手改（改源 MD 后重跑）。

\begin{gknotice}
\textbf{使用说明}：本节只给最终答案便于快速核对，完整推导见解析版。答案为编写者独立推导，\textbf{非官方标准答案}。
\end{gknotice}

\gkpart{选择题}
\begin{center}\begin{tabular}{ccccccccc}
\toprule
题号 & 1 & 2 & 3 & 4 & 5 & 6 & 7 & 8 \\
\midrule
答案 & C & D & B & A & B & A & D & D \\
\bottomrule
\end{tabular}\end{center}

\gkpart{填空题}
\begin{itemize}[leftmargin=2.6em]
  \item[\textbf{9.}] $(-\infty,-3]\cup[2,+\infty)$。按 $x\leqslant-2$、$-2<x<1$、$x\geqslant1$ 分段去绝对值，分别得 $x\leqslant-3$、无解、$x\geqslant2$。
  \item[\textbf{10.}] $y=-5x+3$（即 $5x+y-3=0$）。$y'=-5\mathrm e^{-5x}$，在 $x=0$ 处斜率为 $-5$，过点 $(0,3)$。
  \item[\textbf{11.}] $\dfrac16$。中位数为 $6$ 须恰取 $3$ 个小于 $6$ 的数和 $3$ 个大于 $6$ 的数，有利取法 $\mathrm C_6^3\mathrm C_3^3=20$，总取法 $\mathrm C_{10}^7=120$。
  \item[\textbf{12.}] $2$。由投影公式 $b\cos C+c\cos B=a$，得 $a=2b$，故 $\dfrac ab=2$。
  \item[\textbf{13.}] $50$。由 $a_9a_{12}=a_{10}a_{11}$ 得 $a_{10}a_{11}=\mathrm e^5$，而 $\prod_{i=1}^{20}a_i=(a_{10}a_{11})^{10}=\mathrm e^{50}$，取对数得 $50$。
  \item[\textbf{14.}] $(1,1)$。$C_1$ 化为 $y^2=x$，$C_2$ 化为 $y=1$，联立得 $x=1$。
  \item[\textbf{15.}] $9$。$\triangle AEF\sim\triangle CDF$，相似比 $\dfrac{AE}{DC}=\dfrac13$，面积比为相似比的平方的倒数，即 $9$。
\end{itemize}

\gkpart{解答题}
\begin{itemize}[leftmargin=2.6em]
  \item[\textbf{16.}] (1) $\sqrt3$。$\sin\left(\dfrac{5\pi}{12}+\dfrac\pi4\right)=\sin\dfrac{2\pi}{3}=\dfrac{\sqrt3}{2}$，故 $A\cdot\dfrac{\sqrt3}{2}=\dfrac32$；(2) $\dfrac{\sqrt{30}}4$。$f(\theta)+f(-\theta)=\sqrt6\cos\theta=\dfrac32$ 得 $\cos\theta=\dfrac{\sqrt6}{4}$，$\sin\theta=\dfrac{\sqrt{10}}4$，所求为 $\sqrt3\sin\theta$
  \item[\textbf{17.}] (1) $n_1=6$，$n_2=3$，$f_1=0.24$，$f_2=0.12$。$(40,45]$ 内有 $6$ 个数据，$(45,50]$ 内有 $3$ 个，频率分别为 $\dfrac6{25}$ 与 $\dfrac3{25}$；(2) 纵轴取 $\dfrac{\text{\text{频率}}}{\text{\text{组距}}}$，五个矩形的高依次为 $0.024$，$0.040$，$0.064$，$0.048$，$0.024$。组距为 $5$，各高等于对应频率除以 $5$；(3) $\dfrac{369}{625}$（约 $0.5904$）。落在 $(30,35]$ 的概率为 $\dfrac15$，任取 $4$ 人至少 $1$ 人落在该区间的概率为 $1-\left(\dfrac45\right)^4$
  \item[\textbf{18.}] (1) 见解析版；(2) $\dfrac{2\sqrt{57}}{19}$。设边长 $1$，$PD=\sqrt3$，取 $D$ 为原点建系得 $F\left(0,\dfrac34,\dfrac{\sqrt3}{4}\right)$，$E\left(0,0,\dfrac{\sqrt3}{4}\right)$，两平面法向量为 $(0,-1,\sqrt3)$ 与 $(\sqrt3,0,4)$
  \item[\textbf{19.}] (1) $a_1=3$，$a_2=5$，$a_3=7$。由 $n=1,2$ 及 $S_3=15$ 联立解得；(2) $a_n=2n+1$。由递推式得 $S_{n+1}=\dfrac{(2n+1)S_n+3n^2+4n}{2n}$，归纳得 $S_n=n^2+2n$，相减得 $a_n=2n+1$
  \item[\textbf{20.}] (1) $\dfrac{x^2}{9}+\dfrac{y^2}{4}=1$。$c=\sqrt5$，由 $\dfrac ca=\dfrac{\sqrt5}{3}$ 得 $a=3$，$b^2=9-5=4$；(2) $x^2+y^2=13$。设切线 $y=kx+m$，相切条件化为 $(x_0^2-9)k^2-2x_0y_0k+(y_0^2-4)=0$，由 $k_1k_2=-1$ 得 $x_0^2+y_0^2=13$（竖直切线情形同样满足）
  \item[\textbf{21.}] (1) $D=(-\infty,-1-\sqrt{2-k})\cup(-1-\sqrt{-2-k},-1+\sqrt{-2-k})\cup(-1+\sqrt{2-k},+\infty)$。令 $t=x^2+2x+k$，由 $t^2+2t-3>0$ 得 $t<-3$ 或 $t>1$，再解出 $x$；(2) 增区间为 $(-\infty,-1-\sqrt{2-k})$ 与 $(-1,-1+\sqrt{-2-k})$；减区间为 $(-1-\sqrt{-2-k},-1)$ 与 $(-1+\sqrt{2-k},+\infty)$。由 $f'(x)=-\dfrac{2(t+1)(x+1)}{(t^2+2t-3)^{3/2}}$ 定号；(3) $\left(-1-\sqrt{-2k-4},-1-\sqrt{-k+2}\right)\cup\left(-1-\sqrt{-k-2},-3\right)\cup\left(1,-1+\sqrt{-k-2}\right)\cup\left(-1+\sqrt{-k+2},-1+\sqrt{-2k-4}\right)$。由 $f(x)>f(1)$ 化为 $|(x+1)^2+k|<-k-4$，再与定义域条件 $|(x+1)^2+k|>2$ 取交
\end{itemize}
